% Chapter Template

\chapter{Conclusiones}

\label{Chapter5}

En este capítulo se presentan las conclusiones del trabajo realizado, el grado
de cumplimiento de los requerimientos, las desviaciones respecto de la
planificación original, los riesgos que se manifestaron durante el desarrollo y
las lecciones aprendidas sobre las técnicas empleadas.

%----------------------------------------------------------------------------------------
%	SECTION 1
%----------------------------------------------------------------------------------------

\section{Conclusiones generales}

El objetivo principal del trabajo fue el desarrollo de un sistema integral de
análisis y optimización biomecánica para ciclistas, lo que integra visión por
computadora, modelos biomecánicos y algoritmos evolutivos multiobjetivo. Los
principales logros obtenidos fueron los siguientes:

\begin{itemize}
    \item Se desarrolló un pipeline completo de adquisición, preprocesamiento y análisis
          de datos biomecánicos a partir de video, con registro temporal sincronizado con
          las métricas del rodillo inteligente.
    \item Se resolvió el problema de \textit{motion blur} mediante la sustitución de la
          cámara original por un dispositivo Google Pixel 8~Pro, con mejoras en la
          iluminación y la estabilidad del soporte.
    \item Se aplicó la estimación de pose con MediaPipe BlazePose sobre más de 26.000
          fotogramas, con superposición visual de los puntos clave para la verificación
          del ingeniero.
    \item Se incorporó un filtro de Kalman univariado para la reducción del ruido en las
          series temporales de los puntos clave, con preservación de la información
          biomecánica relevante.
    \item Se construyó un gemelo digital del ciclista con cinemática inversa, lo que
          permitió simular el impacto de distintas configuraciones de la bicicleta sobre
          las métricas biomecánicas.
    \item Se implementó el algoritmo NSGA-II para la optimización multiobjetivo de la
          posición del ciclista, con balance entre confort, potencia y aerodinámica sobre
          el frente de Pareto.
    \item Se incorporó la morfología del ciclista como variable de entrada en el modelo
          de optimización, lo que permitió la generación de recomendaciones
          personalizadas.
    \item Se desarrolló una interfaz web lo que guía al usuario en todo el proceso, desde
          la captura de video hasta la descarga del reporte de recomendaciones.
    \item Se alcanzó un grado de cumplimiento de los requerimientos del 89\,\% (26 de 29
          criterios de aceptación cumplidos), con un criterio parcialmente cumplido y dos
          no realizados, correspondientes a la suite de pruebas automatizadas.
\end{itemize}

La planificación original se siguió de manera satisfactoria hasta la etapa de
selección de cámara, donde surgieron inconvenientes lo que retrasaron el
cronograma. Durante esa etapa se manifestaron dos de los riesgos identificados
en la planificación: la baja calidad de imagen por efecto de \textit{motion
    blur} y la consecuente inexactitud del modelo de pose. La cámara seleccionada
inicialmente no fue diseñada para la captura a alta velocidad, lo que derivó en
imágenes borrosas no detectables por el modelo. El plan de mitigación resultó
efectivo: se reemplazó el dispositivo de captura por el Google Pixel 8~Pro, se
mejoró la iluminación del entorno y se aumentó la estabilidad del soporte, lo
que permitió retomar el cronograma con un retraso acotado.

En cuanto a las técnicas empleadas, la metodología CRISP-DM resultó de utilidad
para la organización y el avance iterativo del proyecto, ya que proveyó una
guía clara para cada etapa del desarrollo. El modelo de pose HRNet, utilizado
en la etapa inicial, no se adaptó a las condiciones de captura del entorno, por
lo que fue reemplazado por MediaPipe BlazePose, lo que demostró mayor robustez
bajo las condiciones reales de trabajo.

%----------------------------------------------------------------------------------------
%	SECTION 2
%----------------------------------------------------------------------------------------

\section{Próximos pasos}

A partir de los resultados obtenidos, se identifican las siguientes líneas de
trabajo para la continuación del proyecto:

\begin{itemize}
    \item La incorporación de una suite de pruebas automatizadas (unitarias, de
          integración y de extremo a extremo) para garantizar la mantenibilidad y la
          calidad del sistema a largo plazo.
    \item La sustitución del dispositivo de captura por una cámara de mayor precisión y
          frecuencia de muestreo, lo que permitiría obtener series temporales más densas
          y mejorar la calidad del análisis biomecánico.
    \item La extensión del modelo biomecánico para la inclusión de las extremidades
          superiores y el análisis del tren superior del ciclista.
    \item La validación del sistema con una muestra ampliada de ciclistas de distintos
          niveles y morfologías, con el fin de evaluar la generalización de las
          recomendaciones.
    \item La integración de sensores adicionales, como plataformas de fuerza en los
          pedales, para enriquecer el modelo biomecánico con datos de carga real.
    \item El despliegue del sistema en un entorno de producción accesible para
          profesionales del ciclismo, con soporte para múltiples usuarios y sesiones
          concurrentes.
\end{itemize}
